% GENERATED FILE. Edit canonical lessons, then run npm run generate:books.
% canonical-source-hash: fnv1a64:0aba490d894a9435
% canonical-lessons: PT-C358-funcionario, PT-C358-turno, PT-C358-sindicato, PT-C358-greve, PT-C358-reivindicar, PT-R358-revisao-351-358

\chapter{A greve}
\label{ch:pt-a-greve}
\hlchaptermodality{358}

\begin{chapteropening}
\textbf{By the end of this chapter:} \emph{I can talk about strikes and unions: staff, shifts and what workers demand.}

Complete the last lesson of chapter 358: I can talk about strikes and unions: staff, shifts and what workers demand.
\end{chapteropening}

\section[o funcionário]{o funcionário — the employee, the official, the clerk}

\label{lesson:PT-C358-funcionario}



\begin{quote}
Before the new one: say the Portuguese for to employ, then the Portuguese for the employer.
\end{quote}

\subsection*{You'll want to know: o funcionário}
\textbf{o funcionário} — \textquotedblleft{}the employee, the official, the clerk\textquotedblright{}.

\begin{quote}
O funcionário do banco foi muito simpático.
\end{quote}

\begin{quote}
\emph{The bank clerk was very nice.}
\end{quote}

Feminine \emph{a funcionária}. \emph{Funcionário público} — civil servant.

\begin{cousinweb}
\emph{funcionário} was modelled on French \emph{fonctionnaire} (EtymDB), built on \emph{função} \textquotedblleft{}function\textquotedblright{}, from Latin \emph{fūnctiō} \textquotedblleft{}a performing, an execution\textquotedblright{} (Lewis \& Short).
\end{cousinweb}

\begin{grammarlens}[title={What you've built}]
One more word for telling, explaining and describing at B1.
\end{grammarlens}

\subsection*{Your turn}
\begin{itemize}
\raggedright
  \item \emph{Say it:} \emph{o funcionário}
  \item \emph{Say it:} \emph{o funcionário}, once more
  \item \emph{Say it:} a sentence of your own with \emph{o funcionário}
  \item \emph{Recall:} say the Portuguese for to employ, then the Portuguese for the employer, then say \emph{o funcionário} again
\end{itemize}

\begin{tcolorbox}[breakable,colback=teal!4,colframe=teal!35!black,title={Before you move on}]
What does o funcionário mean? (\textquotedblleft{}The employee, the official, the clerk\textquotedblright{}.) Say it once more.
\end{tcolorbox}
\section[o turno]{o turno — the shift; the turn}

\label{lesson:PT-C358-turno}



\begin{quote}
Before the new one: say the Portuguese for the employer, then the Portuguese for the employee.
\end{quote}

\subsection*{You'll want to know: o turno}
\textbf{o turno} — \textquotedblleft{}the shift; the turn\textquotedblright{}.

\begin{quote}
Esta semana trabalho no turno da noite.
\end{quote}

\begin{quote}
\emph{This week I am on the night shift.}
\end{quote}

\emph{Por turnos} — in shifts.

\begin{cousinweb}
\emph{turno} was borrowed from Spanish \emph{turno} (EtymDB), built on \emph{turnar} \textquotedblleft{}to take turns\textquotedblright{}, which goes back to French \emph{tourner}, from Latin \emph{tornāre} \textquotedblleft{}to turn on a lathe\textquotedblright{} (Lewis \& Short).
\end{cousinweb}

\begin{grammarlens}[title={What you've built}]
One more word for telling, explaining and describing at B1.
\end{grammarlens}

\subsection*{Your turn}
\begin{itemize}
\raggedright
  \item \emph{Say it:} \emph{o turno}
  \item \emph{Say it:} \emph{o turno}, once more
  \item \emph{Say it:} a sentence of your own with \emph{o turno}
  \item \emph{Recall:} say the Portuguese for the employer, then the Portuguese for the employee, then say \emph{o turno} again
\end{itemize}

\begin{tcolorbox}[breakable,colback=teal!4,colframe=teal!35!black,title={Before you move on}]
What does o turno mean? (\textquotedblleft{}The shift; the turn\textquotedblright{}.) Say it once more.
\end{tcolorbox}
\section[o sindicato]{o sindicato — the trade union}

\label{lesson:PT-C358-sindicato}



\begin{quote}
Before the new one: say the Portuguese for the employee, then the Portuguese for the shift.
\end{quote}

\subsection*{You'll want to know: o sindicato}
\textbf{o sindicato} — \textquotedblleft{}the trade union\textquotedblright{}.

\begin{quote}
Ela entrou para o sindicato no primeiro ano.
\end{quote}

\begin{quote}
\emph{She joined the union in her first year.}
\end{quote}

\emph{O sindicato dos professores} — the teachers' union.

\begin{cousinweb}
\emph{sindicato} is built on \emph{síndico} \textquotedblleft{}representative\textquotedblright{}, from Latin \emph{syndicus}, from Greek \emph{súndikos} \textquotedblleft{}advocate\textquotedblright{}; French \emph{syndicat} \textquotedblleft{}trade union\textquotedblright{}, from Medieval Latin \emph{syndicātus}, shaped the modern sense (EtymDB).
\end{cousinweb}

\begin{grammarlens}[title={What you've built}]
One more word for telling, explaining and describing at B1.
\end{grammarlens}

\subsection*{Your turn}
\begin{itemize}
\raggedright
  \item \emph{Say it:} \emph{o sindicato}
  \item \emph{Say it:} \emph{o sindicato}, once more
  \item \emph{Say it:} a sentence of your own with \emph{o sindicato}
  \item \emph{Recall:} say the Portuguese for the employee, then the Portuguese for the shift, then say \emph{o sindicato} again
\end{itemize}

\begin{tcolorbox}[breakable,colback=teal!4,colframe=teal!35!black,title={Before you move on}]
What does o sindicato mean? (\textquotedblleft{}The trade union\textquotedblright{}.) Say it once more.
\end{tcolorbox}
\section[a greve]{a greve — the strike}

\label{lesson:PT-C358-greve}



\begin{quote}
Before the new one: say the Portuguese for the shift, then the Portuguese for the trade union.
\end{quote}

\subsection*{You'll want to know: a greve}
\textbf{a greve} — \textquotedblleft{}the strike\textquotedblright{}.

\begin{quote}
Os comboios não circulam hoje por causa da greve.
\end{quote}

\begin{quote}
\emph{The trains are not running today because of the strike.}
\end{quote}

\emph{Fazer greve} — to go on strike; \emph{estar em greve} — to be on strike.

\begin{cousinweb}
\emph{greve} was borrowed from French \emph{grève}, which first meant \textquotedblleft{}a sandy riverbank\textquotedblright{}; the sense \textquotedblleft{}strike\textquotedblright{} is usually explained by a riverside square in Paris where men out of work gathered. EtymDB lists the Portuguese word but gives no source for it.
\end{cousinweb}

\begin{grammarlens}[title={What you've built}]
One more word for telling, explaining and describing at B1.
\end{grammarlens}

\subsection*{Your turn}
\begin{itemize}
\raggedright
  \item \emph{Say it:} \emph{a greve}
  \item \emph{Say it:} \emph{a greve}, once more
  \item \emph{Say it:} a sentence of your own with \emph{a greve}
  \item \emph{Recall:} say the Portuguese for the shift, then the Portuguese for the trade union, then say \emph{a greve} again
\end{itemize}

\begin{tcolorbox}[breakable,colback=teal!4,colframe=teal!35!black,title={Before you move on}]
What does a greve mean? (\textquotedblleft{}The strike\textquotedblright{}.) Say it once more.
\end{tcolorbox}
\section[reivindicar]{reivindicar — to lay claim to, to demand (as a right)}

\label{lesson:PT-C358-reivindicar}



\begin{quote}
Before the new one: say the Portuguese for the trade union, then the Portuguese for the strike.
\end{quote}

\subsection*{You'll want to know: reivindicar}
\textbf{reivindicar} — \textquotedblleft{}to lay claim to, to demand (as a right)\textquotedblright{}.

\begin{quote}
Os trabalhadores reivindicam salários mais altos.
\end{quote}

\begin{quote}
\emph{The workers are demanding higher wages.}
\end{quote}

Formal, of rights and pay; in speech \emph{pedir}.

\begin{cousinweb}
\emph{reivindicar} is formed from Latin \emph{rēs, reī} \textquotedblleft{}thing\textquotedblright{} and \emph{vindicāre} \textquotedblleft{}to lay claim to, to avenge\textquotedblright{} (EtymDB, Lewis \& Short): to claim a thing as yours.
\end{cousinweb}

\begin{grammarlens}[title={What you've built}]
One more word for telling, explaining and describing at B1.
\end{grammarlens}

\subsection*{Your turn}
\begin{itemize}
\raggedright
  \item \emph{Say it:} \emph{reivindicar}
  \item \emph{Say it:} \emph{reivindicar}, once more
  \item \emph{Say it:} a sentence of your own with \emph{reivindicar}
  \item \emph{Recall:} say the Portuguese for the trade union, then the Portuguese for the strike, then say \emph{reivindicar} again
\end{itemize}

\begin{tcolorbox}[breakable,colback=teal!4,colframe=teal!35!black,title={Before you move on}]
What does reivindicar mean? (\textquotedblleft{}To lay claim to, to demand (as a right)\textquotedblright{}.) Say it once more.
\end{tcolorbox}
\section[(dialogue)]{Revisão dos capítulos 351–358}

\label{lesson:PT-R358-revisao-351-358}



\begin{quote}
Say the Portuguese for the impact and the repercussion.
\end{quote}

\subsection*{Your turn}
Say these aloud:

\begin{itemize}
\raggedright
  \item the impact, the repercussion, to involve, to entail, to trigger
  \item the reaction, to react, the pretext, the excuse, to justify
  \item the obstacle, the challenge, to prevent, to make difficult, to favour
  \item the influence, to influence, to make worse, to ease, to make up for
  \item the vacancy, the CV, the application, the candidate, to hire
  \item to fire, the resignation, the dismissal, unemployed, the unemployment
  \item the internship, the intern, the apprentice, to employ, the employer
  \item the employee, the shift, the trade union, the strike, to lay claim to
\end{itemize}

\begin{tcolorbox}[breakable,colback=teal!4,colframe=teal!35!black,title={Before you move on}]
The impact? (\textbf{o impacto}.) The repercussion? (\textbf{a repercussão}.) To involve? (\textbf{implicar}.) To entail? (\textbf{acarretar}.) To trigger? (\textbf{desencadear}.) The reaction? (\textbf{a reação}.) To react? (\textbf{reagir}.) The pretext? (\textbf{o pretexto}.) The excuse? (\textbf{a desculpa}.) To justify? (\textbf{justificar}.) The obstacle? (\textbf{o obstáculo}.) The challenge? (\textbf{o desafio}.) To prevent? (\textbf{impedir}.) To make difficult? (\textbf{dificultar}.) To favour? (\textbf{favorecer}.) The influence? (\textbf{a influência}.) To influence? (\textbf{influenciar}.) To make worse? (\textbf{agravar}.) To ease? (\textbf{atenuar}.) To make up for? (\textbf{compensar}.) The vacancy? (\textbf{a vaga}.) The CV? (\textbf{o currículo}.) The application? (\textbf{a candidatura}.) The candidate? (\textbf{o candidato}.) To hire? (\textbf{contratar}.) To fire? (\textbf{despedir}.) The resignation? (\textbf{a demissão}.) The dismissal? (\textbf{o despedimento}.) Unemployed? (\textbf{desempregado}.) The unemployment? (\textbf{o desemprego}.) The internship? (\textbf{o estágio}.) The intern? (\textbf{o estagiário}.) The apprentice? (\textbf{o aprendiz}.) To employ? (\textbf{empregar}.) The employer? (\textbf{a entidade patronal}.) The employee? (\textbf{o funcionário}.) The shift? (\textbf{o turno}.) The trade union? (\textbf{o sindicato}.) The strike? (\textbf{a greve}.) To lay claim to? (\textbf{reivindicar}.)
\end{tcolorbox}
